\documentclass[11pt]{article}

\usepackage[T1]{fontenc}
\usepackage[utf8]{inputenc}
\usepackage{lmodern}
\usepackage{amsmath,amssymb,amsthm}
\usepackage{booktabs}
\usepackage{enumitem}
\usepackage[margin=1in]{geometry}
\usepackage{microtype}
\usepackage{xcolor}
\usepackage{xurl}
\usepackage[colorlinks=true,linkcolor=blue!55!black,citecolor=blue!55!black,
  urlcolor=blue!55!black]{hyperref}
\hypersetup{
  pdftitle={Exact Local Exchange Profiles and Reconfiguration Barriers for Isosceles-Triangle-Free Grid Sets},
  pdfauthor={Pierre-Baptiste Borges},
  pdfsubject={Local exchange profiles, TAR barriers, and machine-checked certificates}
}

\newtheorem{theorem}{Theorem}
\newtheorem{lemma}[theorem]{Lemma}
\newtheorem{proposition}[theorem]{Proposition}
\newtheorem{corollary}[theorem]{Corollary}
\newtheorem{computationalclaim}[theorem]{Computational claim}
\theoremstyle{definition}
\newtheorem{definition}[theorem]{Definition}
\theoremstyle{remark}
\newtheorem{remark}[theorem]{Remark}

\newcommand{\F}{\mathcal{F}}
\newcommand{\calU}{\mathcal{U}}
\newcommand{\calC}{\mathcal{C}}
\newcommand{\calM}{\mathcal{M}}
\newcommand{\dist}{d_{\triangle}}
\newcommand{\N}{\mathbb{N}}

\title{Exact Local Exchange Profiles and Reconfiguration Barriers for
Isosceles-Triangle-Free Grid Sets}
\author{Pierre-Baptiste Borges\\[2pt]\small Independent researcher\\
  \small\texttt{pierrebaptiste.borges@gmail.com}}
\date{27 August 2026}

\begin{document}

\maketitle

\begin{abstract}
Let $S$ be a fixed set of points in an integer grid containing no isosceles
triangle, with collinear triples included.  We introduce two local invariants.
The individual-unlock profile $u_S(r)$ is the largest number of outside points
that become individually admissible after deleting $r$ points of $S$.  The
exchange profile $e_S(r)$ is the largest number of mutually compatible outside
points that can be inserted after those deletions.  We prove an exact slice
identity for $e_S$ and a general lower bound on token-addition/removal (TAR)
reconfiguration barriers derived from either profile.  We then describe an
exact algorithm based on link-graph vertex covers, pair-forced deletion masks,
and antichains of minimal masks.

For the public $112$-point configuration in the $64\times64$ grid, our current
computation gives $e_S(7)=5<u_S(7)=6$, showing that individual admissibility can
strictly overestimate repair capacity.  For the public $164$-point
$100\times100$ configuration, an independently generated CNF and a verified
DRAT proof establish $u_S(6)=4$.  At radius seven, a directly checked witness
gives $e_S(7)\ge5$, while an independently regenerated core-SAT instance and a
second verified DRAT proof establish $e_S(7)<6$.  At radius eight, a directly
checked $8\to6$ exchange and an independently audited exhaustive search give
$e_S(8)=6$; the negative search remains program-verified rather than
proof-trace certified.  It follows that every distinct isosceles-free set of
size at least $164$ deletes at least nine points of the $164$-point
configuration, every equal-size alternative has symmetric-difference distance
at least $18$, every improvement has distance at least $19$, and every TAR
path to an equal-or-larger distinct set descends to size at most $161$.  We
make no claim that either public configuration is globally optimal.
\end{abstract}

\begin{center}
\fbox{\parbox{0.92\linewidth}{\textbf{Computational evidence.}
The radius-six and radius-seven upper bounds for the
$100\times100$ instance have machine-verified DRAT certificates and
independently regenerated reductions.  The new radius-eight upper bound is an
independently audited exhaustive program computation, but has no DRAT/LRAT
proof trace.  The results therefore have different evidence levels;
Section~\ref{sec:certification} states those levels and the remaining trust
boundaries.}}
\end{center}

\section{Introduction}

Write
\[
  G_n=\{0,1,\ldots,n-1\}^2.
\]
The extremal problem asks for a largest subset of $G_n$ containing no three
distinct points whose three pairwise distances include a repeated value.
Thus an isosceles triangle is forbidden even when its three vertices are
collinear.  The problem has recently served as a benchmark for machine-assisted
construction.  PatternBoost reported SAT optima for small grids, a
$110$-point set for $n=64$, and a $160$-point set for $n=100$
\cite{charton2024patternboost}.  AlphaEvolve subsequently published the
$112$- and $164$-point configurations studied here and expressly did not claim
that the latter is optimal \cite{georgiev2025exploration,alphaevolve59}.
Geometry-aware Monte Carlo tree search has also treated this problem as one of
several combinatorial-geometry benchmarks \cite{zhang2026geometry}.

The global and asymptotic questions remain difficult.  Croot, Mao, Pohoata,
Sheffer, and Yip proved an upper bound of the form
\[
  n^2\exp\!\left(-c\frac{\log n}{\log\log n}\right)
\]
for the maximum cardinality, for an absolute constant $c>0$
\cite{croot2026largesieve}.  Counting all isosceles triangles in a regular grid
is a related but distinct problem \cite{wu2016counting}.  Our focus is neither
global optimality nor asymptotics.  We ask instead how a particular large
construction can be changed locally.

Local optimality and bounded swaps are standard concepts in combinatorial
optimization \cite{komusiewicz2025local,ganian2026framework}.  TAR
reconfiguration, buffer size, cardinality range, and minimum-cardinality paths
are also established notions
\cite{deberg2016thresholds,bulteau2021parametrization,hoang2026maxmin}.
Our candidate contribution is a concrete local profile around a fixed
hypergraph independent set, its translation into a TAR-barrier lower bound,
and exact computations for the two public grid configurations.  Link graphs,
which underlie our reduction, are standard hypergraph objects; see, for
example, \cite{hancock2016solutionfree}.

The contributions of this article are as follows.
\begin{enumerate}[leftmargin=*,itemsep=2pt]
  \item We define the individual-unlock profile $u_S$ and the compatible
  exchange profile $e_S$, and identify $e_S$ with the exact maximum on a
  fixed deletion slice.
  \item We derive a transition-sensitive profile-to-TAR lemma.  The resulting
  expression is a lower bound on the true reconfiguration barrier, not an
  assertion that the barrier itself is exact.
  \item We give an exact antichain recurrence for deciding whether $k$ outside
  points can coexist after at most $r$ deletions in any $3$-uniform
  hypergraph.
  \item We report the current profiles for the AlphaEvolve $n=64$ and $n=100$
  constructions, including the strict separation $e_S(7)<u_S(7)$ at $n=64$
  and the exact value $e_{S_{100}}(8)=6$.
  \item We archive verified DRAT proofs for the independently audited
  individual-unlock cell $u_{S_{100}}(6)=4$ and for the decisive compatible
  upper bound $e_{S_{100}}(7)<6$.  We separately document and audit the
  exhaustive radius-eight computation, which does not yet have the same
  proof-trace evidence level.
\end{enumerate}

To our knowledge after a targeted literature search, these exact profiles and
the particular profile-derived expression in Lemma~\ref{lem:tar} have not been
used for this grid problem.  This is a provisional priority statement, not a
claim that the constituent ideas---local exchange, link graphs, TAR, or SAT
certification---are new.

\section{The hypergraph model and local profiles}
\label{sec:definitions}

Let $H_n=(G_n,E_n)$ be the $3$-uniform hypergraph in which
$\{a,b,c\}\in E_n$ if and only if $a,b,c$ are distinct and
\[
 \|a-b\|_2^2=\|a-c\|_2^2,
 \quad\text{or}\quad
 \|a-b\|_2^2=\|b-c\|_2^2,
 \quad\text{or}\quad
 \|a-c\|_2^2=\|b-c\|_2^2.
\]
All computations use these integer equalities; no floating-point geometry is
needed.  Let $\F_n$ be the hereditary family of independent sets of $H_n$.
Most of the statements below hold for an arbitrary finite ground set $V$ and
an arbitrary hereditary family $\F\subseteq 2^V$.

Fix $S\in\F$, write $m=|S|$, and let $R\subseteq S$ be a set of deleted
original vertices.  Define the set of individually admissible outsiders by
\[
  \calU_S(R)=\bigl\{p\in V\setminus S:
       (S\setminus R)\cup\{p\}\in\F\bigr\}.
\]

\begin{definition}[Individual-unlock and exchange profiles]
For $0\le r\le m$, define
\begin{align*}
  u_S(r)&=\max_{\substack{R\subseteq S\\|R|=r}}|\calU_S(R)|,\\
  e_S(r)&=\max\bigl\{|A|: R\subseteq S,\ |R|=r,\ A\subseteq V\setminus S,
                 \ (S\setminus R)\cup A\in\F\bigr\}.
\end{align*}
Then $e_S(r)\le u_S(r)$.
\end{definition}

The inequality can be strict because two points can each be admissible with
the retained portion of $S$ but form a forbidden triple together with a third
point, or three outsiders can themselves form a forbidden triple.  The
$n=64$, $r=7$ computation in Section~\ref{sec:results} gives a concrete strict
example.

For $X\in\F$, put $R_X=S\setminus X$, $A_X=X\setminus S$, and
$r_X=|R_X|$.  Heredity implies
\[
  A_X\subseteq\calU_S(R_X),
\]
since $(S\setminus R_X)\cup\{p\}\subseteq X$ for each $p\in A_X$.
Consequently
\begin{equation}
  |X|=m-r_X+|A_X|\le m-r_X+u_S(r_X).
  \label{eq:unlock-slice}
\end{equation}
The compatible profile gives the exact maximum rather than only an upper
bound.

\begin{proposition}[Exact slice identity]
\label{prop:slice}
For every $0\le r\le m$,
\[
 \max_{\substack{X\in\F\\|S\setminus X|=r}}|X|
 =m-r+e_S(r).
\]
\end{proposition}

\begin{proof}
Every set in the slice decomposes uniquely as
$X=(S\setminus R_X)\cup A_X$, with $|R_X|=r$ and
$A_X\subseteq V\setminus S$.  Maximizing $|X|=m-r+|A_X|$ over exactly these
feasible pairs $(R_X,A_X)$ is the definition of $e_S(r)$.
\end{proof}

We use symmetric-difference distance
\[
  \dist(X,Y)=|X\triangle Y|.
\]
In this paper, ``$k$-swap-optimal for improvement'' means that no
$T\in\F$ with $|T|>|S|$ satisfies $\dist(S,T)\le k$.  This convention should
be stated explicitly: other local-search papers may bound the numbers deleted
and inserted separately, in which case the numerical radius must be
translated rather than copied.

\section{A profile-derived TAR barrier}
\label{sec:tar}

A TAR path is a sequence $X_0,X_1,\ldots,X_\ell\in\F$ satisfying
$|X_i\triangle X_{i+1}|=1$ for every $i$.  For an integer $R\ge1$, define
\begin{equation}
\begin{aligned}
 D^u_S(R)&=\max\left\{0,\max_{0\le r<R}\bigl(r+1-u_S(r)\bigr)\right\},\\
 D^e_S(R)&=\max\left\{0,\max_{0\le r<R}\bigl(r+1-e_S(r)\bigr)\right\}.
\end{aligned}
\label{eq:barriers}
\end{equation}
Because $e_S(r)\le u_S(r)$, one has $D^e_S(R)\ge D^u_S(R)$.

\begin{lemma}[Deletion-transition TAR barrier]
\label{lem:tar}
Let $V$ be finite, let $\F\subseteq2^V$ be hereditary, and let
$S\in\F$ have size $m$.  Suppose
$X_0=S,X_1,\ldots,X_\ell=T$ is a TAR path and
$|S\setminus T|\ge R$.  Then
\begin{align}
 \min_{0\le i\le\ell}|X_i|
   &\le m-D^u_S(R), \label{eq:tar-u}\\
 \min_{0\le i\le\ell}|X_i|
   &\le m-D^e_S(R). \label{eq:tar-e}
\end{align}
Thus $D^u_S(R)$ and $D^e_S(R)$ are lower bounds on the loss below $|S|$
forced along every such path.
\end{lemma}

\begin{proof}
Set $r_i=|S\setminus X_i|$.  We have $r_0=0$ and $r_\ell\ge R$.
A one-vertex TAR move changes $r_i$ by at most one.  Therefore, for every
$r=0,1,\ldots,R-1$, the path contains an upward transition with
$r_i=r$ and $r_{i+1}=r+1$.  Such a transition removes one retained element of
$S$, and hence $|X_{i+1}|=|X_i|-1$.

Immediately before this deletion, inequality~\eqref{eq:unlock-slice} gives
\[
 |X_i|\le m-r+u_S(r)=m-\bigl(r-u_S(r)\bigr).
\]
Consequently
\[
 |X_{i+1}|\le m-r-1+u_S(r)
 =m-\bigl(r+1-u_S(r)\bigr).
\]
Taking the most restrictive of these mandatory deletion transitions, together
with the trivial $\min_i|X_i|\le|X_0|=m$, yields \eqref{eq:tar-u}.
Proposition~\ref{prop:slice} replaces the pre-deletion slice bound by its exact
maximum $m-r+e_S(r)$ and gives \eqref{eq:tar-e}.
\end{proof}

\begin{remark}[What the formula does and does not prove]
Even when every value of $u_S$ or $e_S$ is exact, $D^u_S(R)$ or
$D^e_S(R)$ is only a lower bound on the true TAR barrier.  Compatibility
between slices and the order of valid moves can force a deeper valley.  If one
has only certified upper bounds $\overline u(r)\ge u_S(r)$, then
$\max\{0,\max_{r<R}(r+1-\overline u(r))\}$ is still a valid, possibly weaker,
lower bound.
The strict range $r<R$ in \eqref{eq:barriers} is essential: a path whose target
deletes at least $R$ originals is guaranteed to make the upward transitions
$r\to r+1$ for $0\le r<R$ without any further hypothesis.  A path reaching
$R+1$ makes one additional transition, so a known radius-$R$ profile value can
strengthen the maximum; we use $r<R$ deliberately because it requires only the
strict prefix.  The maximum with zero in \eqref{eq:barriers} also covers a
nonmaximal $S$, for which additions before the first deletion may prevent any
loss below $m$.  The configurations studied here have $u_S(0)=0$.
\end{remark}

\begin{corollary}[Local isolation and improvement distance]
\label{cor:distance}
Assume $e_S(0)=0$, and let $R\ge1$.
\begin{enumerate}[label=(\roman*),leftmargin=*]
  \item If $e_S(r)<r$ for $1\le r<R$, then every $T\in\F\setminus\{S\}$
  with $|T|\ge|S|$ deletes at least $R$ elements of $S$.
  \item If $e_S(r)\le r$ for $0\le r<R$, then every strict improvement
  $T\in\F$ with $|T|>|S|$ deletes at least $R$ elements of $S$ and satisfies
  $\dist(S,T)\ge2R+1$.
  \item Under the hypothesis in (i), every distinct equal-size $T$ satisfies
  $\dist(S,T)\ge2R$.
\end{enumerate}
\end{corollary}

\begin{proof}
Let $r=|S\setminus T|$ and $a=|T\setminus S|$.  Feasibility gives
$a\le e_S(r)$.  If $|T|\ge|S|$, then $a\ge r$, contradicting
$e_S(r)<r$ when $1\le r<R$; the $r=0$ case has $T=S$ because $e_S(0)=0$.
For a strict improvement, $a\ge r+1$; hence
$e_S(r)\le r$ excludes $r<R$, and
$\dist(S,T)=r+a\ge2R+1$.  For a distinct equal-size set, $a=r$ and the same
identity gives distance at least $2R$.
\end{proof}

\section{Exact exchange computation by antichains}
\label{sec:algorithm}

This section gives the mathematical recurrence used by the radius-seven and
radius-eight searches.  It is independent of grid geometry and applies to a fixed independent
set $S$ in any $3$-uniform hypergraph $H=(V,E)$.

\subsection{Link graphs and individual covers}

For $p\in V\setminus S$, define its link graph on $S$ by
\[
  L_S(p)=\bigl(S,\{\{a,b\}\subseteq S:\{p,a,b\}\in E\}\bigr).
\]
Then
\begin{equation}
 p\in\calU_S(R)
 \quad\Longleftrightarrow\quad
 R\text{ is a vertex cover of }L_S(p).
 \label{eq:link-cover}
\end{equation}
Indeed, the only possible hyperedge in $(S\setminus R)\cup\{p\}$ contains
$p$ and two retained vertices of $S$.

Fix a deletion radius $r$.  Let $\calC_r(p)$ be the family of
inclusion-minimal vertex covers of $L_S(p)$ having size at most $r$.  A
branching recursion that selects an uncovered edge $ab$ and branches on
deleting $a$ or $b$ enumerates this family exactly.  Lower bounds from
edge-disjoint matchings can prune a branch without changing the answer.

\subsection{Pair-forced masks and the antichain recurrence}

For two outsiders $p,q\in V\setminus S$, put
\[
  F_S(p,q)=\{s\in S:\{p,q,s\}\in E\}.
\]
Every deletion set supporting both $p$ and $q$ must contain the entire mask
$F_S(p,q)$.

For a family $\mathcal A\subseteq2^S$, write
$\min_{\subseteq}(\mathcal A)$ for its inclusion-minimal members.  Define
$\calM_r(\varnothing)=\{\varnothing\}$.  If $Q\subseteq V\setminus S$ has no
hyperedge entirely among its elements and $p\notin Q$, define
\begin{multline}
 \calM_r(Q\cup\{p\})=\min_{\subseteq}\Bigl\{
 M\cup C\cup\bigcup_{q\in Q}F_S(p,q):
 M\in\calM_r(Q),\ C\in\calC_r(p),\\
 \bigl|M\cup C\cup\bigcup_{q\in Q}F_S(p,q)\bigr|\le r
 \Bigr\}.
 \label{eq:antichain-recurrence}
\end{multline}
If $Q\cup\{p\}$ contains an all-outsider hyperedge, set
$\calM_r(Q\cup\{p\})=\varnothing$.

\begin{proposition}[Exactness of the antichain recurrence]
\label{prop:antichain}
For every $Q\subseteq V\setminus S$, $\calM_r(Q)$ is precisely the antichain
of inclusion-minimal sets $M\subseteq S$ of size at most $r$ such that
$(S\setminus M)\cup Q$ is independent.  Consequently,
\[
 e_S(r)\ge k
 \quad\Longleftrightarrow\quad
 \text{there exists }Q\subseteq V\setminus S, |Q|=k,
 \ \calM_r(Q)\ne\varnothing.
\]
Here a supporting mask of size less than $r$ may be padded with arbitrary
additional deletions from $S$.
\end{proposition}

\begin{proof}
The claim is immediate for $Q=\varnothing$.  Suppose it holds for $Q$.
Any forbidden triple in $(S\setminus M')\cup(Q\cup\{p\})$ is of one of four
types.  A triple contained in $S$ is impossible because $S$ is independent.
A triple containing $p$ and two vertices of $S$ is excluded exactly when
$M'$ contains a cover in $\calC_r(p)$, by \eqref{eq:link-cover}.  A triple
containing $p,q$ and one vertex $s\in S$ is excluded exactly when
$F_S(p,q)\subseteq M'$.  Conditions involving old elements of $Q$ are already
encoded by a member $M\in\calM_r(Q)$.  Finally, a triple entirely in
$Q\cup\{p\}$ cannot be repaired by deleting elements of $S$ and is rejected
explicitly.

Thus the unions in \eqref{eq:antichain-recurrence} are exactly the supporting
masks of size at most $r$, before removal of supersets.  Taking
inclusion-minimal members preserves existence and yields precisely the claimed
antichain.  The final equivalence follows from the definition of $e_S(r)$ and
the fact that independence is preserved under additional deletions.
\end{proof}

The implementation first constructs a graph on eligible outsiders, retaining
an edge only when its two endpoints have a common supporting mask of size at
most $r$.  A feasible $k$-set must be a $k$-clique in this graph, so a
clique-style DFS provides an efficient necessary prefilter.  Pairwise
compatibility is not sufficient: at every node the DFS carries the full
antichain from \eqref{eq:antichain-recurrence} and checks all-outsider triples.

\section{Computational results}
\label{sec:results}

The two main configurations are the public AlphaEvolve outputs
\cite{georgiev2025exploration,alphaevolve59}.  We denote the $164$-point set in
$G_{100}$ by $S_{100}$ and the $112$-point set in $G_{64}$ by $S_{64}$.
Section~\ref{sec:certification} separates theorem-level deductions from the
current evidence for each computed cell.

\subsection{The \texorpdfstring{$164$}{164}-point construction in
\texorpdfstring{$G_{100}$}{G100}}

The current profile is
\begin{center}
\begin{tabular}{c*{9}{c}}
\toprule
$r$ & $0$ & $1$ & $2$ & $3$ & $4$ & $5$ & $6$ & $7$ & $8$\\
\midrule
$u_{S_{100}}(r)$ & $0$ & $0$ & $1$ & $1$ & $2$ & $3$ & $4$ & $\ge5$ & $\ge6$\\
$e_{S_{100}}(r)$ & $0$ & $0$ & $1$ & $1$ & $2$ & $3$ & $4$ & $5$ & $6$\\
\bottomrule
\end{tabular}
\end{center}
Only the displayed lower bounds are claimed for $u_{S_{100}}(7)$ and
$u_{S_{100}}(8)$; their exact values remain open.  Time-limited decisions at
both radii returned \textsc{unknown}, which is not an upper bound.  There is a
directly checked $7\to5$ exchange:
\begin{align*}
R={}&\{(92,51),(96,97),(7,51),(8,93),(26,40),(73,40),(77,95)\},\\
A={}&\{(95,98),(89,56),(75,81),(10,56),(77,99)\}.
\end{align*}
Exact integer geometry verifies that $(S_{100}\setminus R)\cup A$ is
independent.  The antichain DFS also reports $e_{S_{100}}(7)<6$, with
$15{,}227$ search states and $47.762$ seconds in the logged run.  The decisive
upper bound no longer rests on that traversal alone: a separately encoded
no-symmetry core CNF and its DRAT proof have passed independent regeneration
and proof checking, as detailed in Section~\ref{sec:certification}.

At radius eight, an independently checked exchange is
\begin{align*}
R={}&\{(26,59),(7,48),(73,59),(22,4),\\
    &\qquad (92,48),(77,4),(96,2),(3,2)\},\\
A={}&\{(4,1),(10,43),(22,0),(77,0),(89,43),(95,1)\}.
\end{align*}
Exact integer geometry verifies that the resulting $162$-point set is
independent.  An exhaustive antichain search rejects seven additions even in
a necessary-condition relaxation that omits all constraints on triples of
outsiders and disables color pruning.  Section~\ref{sec:certification}
records the independent input regeneration, search counts, and remaining
program-verification boundary.

\begin{computationalclaim}[Audited local consequences for $S_{100}$]
\label{claim:n100}
The radius-eight computation described in Section~\ref{sec:certification}
implies:
\begin{enumerate}[label=(\roman*),leftmargin=*]
  \item every $T\ne S_{100}$ with $|T|\ge164$ deletes at least nine points of
  $S_{100}$;
  \item every distinct $164$-point set has
  $\dist(S_{100},T)\ge18$;
  \item every set of size at least $165$ has
  $\dist(S_{100},T)\ge19$;
  \item $S_{100}$ is $18$-swap-optimal for improvement under the
  symmetric-difference convention of Section~\ref{sec:definitions}; and
  \item every TAR path from $S_{100}$ to a distinct set of size at least $164$
  contains a set of size at most $161$, so no such endpoint lies in the same
  component when TAR is restricted to sets of size at least $162$;
  \item the component of $S_{100}$ in the TAR graph induced by sets of size at
  least $163$ is exactly the star with center $S_{100}$ and leaves
  $S_{100}\setminus\{s\}$ for $s\in S_{100}$.
\end{enumerate}
\end{computationalclaim}

\begin{proof}
For $1\le r\le8$, the displayed exchange values satisfy
$e_{S_{100}}(r)<r$.  Corollary~\ref{cor:distance} with $R=9$ gives (i)--(iv).
For (v), the endpoint must delete at least nine original points by (i), and
Lemma~\ref{lem:tar} applies.  In fact the already computed slice $r=3$ has
$e_{S_{100}}(3)=u_{S_{100}}(3)=1$.  Immediately after the mandatory
transition from three to four deleted originals, the path therefore has size
at most $164-3+1-1=161$.  Equivalently, the known profile entries give
$D^u_{S_{100}}(9)=D^e_{S_{100}}(9)=3$: the open values satisfy
$u_{S_{100}}(7)\ge5$ and $u_{S_{100}}(8)\ge6$, so their contributions to
$D^u$ cannot exceed $3$.

For (vi), $e_{S_{100}}(0)=0$ means that no outsider can be added directly to
$S_{100}$, so its only TAR neighbors are its $164$ single deletions.  Such a
leaf has size $163$: a further deletion violates the threshold, adding an
outsider is impossible because $e_{S_{100}}(1)=0$, and restoring its unique
missing record point returns to $S_{100}$.
\end{proof}

\paragraph{Exploratory radius-nine frontier.}
The next improvement decision is $e_{S_{100}}(9)\ge10$.  A bounded exact DFS
returned \textsc{unknown} after $580.099$ seconds, $233{,}472$ nodes, and
$91{,}148{,}661{,}627$ cover tests; three deterministic heuristics evaluated
$8{,}446{,}705$ deletion masks without finding a $9\to10$ witness.  These are
timeouts and sampled failures, not an upper bound.  They did produce a directly
verified lower bound $u_{S_{100}}(9)\ge7$.  For one archived mask the seven
unlocked outsiders are pairwise compatible, but
\[
  (4,98),\ (75,81),\ (77,99)
\]
form an isosceles triple: the first point has squared distance $5330$ to each
of the other two.  Hence the maximum compatible subset of this particular
seven-point list is six.  This gives a concrete radius-nine illustration that
pair compatibility alone does not capture all repair obstructions; globally,
both profile values at radius nine remain open beyond their lower bounds.

None of these statements asserts that $164$ is the maximum cardinality in
$G_{100}$.  They describe the neighborhood of this one public construction.

\subsection{The \texorpdfstring{$112$}{112}-point construction in
\texorpdfstring{$G_{64}$}{G64}}

For $S_{64}$, the current exact computation is
\begin{center}
\begin{tabular}{c*{8}{c}}
\toprule
$r$ & $0$ & $1$ & $2$ & $3$ & $4$ & $5$ & $6$ & $7$\\
\midrule
$u_{S_{64}}(r)$ & $0$ & $1$ & $2$ & $2$ & $4$ & $4$ & $5$ & $6$\\
$e_{S_{64}}(r)$ & $0$ & $1$ & $2$ & $2$ & $4$ & $4$ & $5$ & $5$\\
\bottomrule
\end{tabular}
\end{center}
In particular,
\[
  e_{S_{64}}(7)=5<6=u_{S_{64}}(7).
\]
This strict separation demonstrates why the individual-unlock profile cannot
in general be treated as the capacity of a simultaneous repair.  The radius
six compatible optimum is also five: $e_{S_{64}}(6)=u_{S_{64}}(6)=5$.

The configuration is not isolated among equal-size sets.  The exchange
\[
  S_{64}'=\bigl(S_{64}\setminus\{(56,2)\}\bigr)\cup\{(62,2)\}
\]
is valid, so $\dist(S_{64},S_{64}')=2$.  On the other hand, any strict
improvement must delete at least eight original points and therefore has
symmetric-difference distance at least $17$.  Thus $S_{64}$ is
$16$-swap-optimal for improvement under our convention.  Moreover,
\[
 D^u_{S_{64}}(8)=2,
 \qquad
 D^e_{S_{64}}(8)=3.
\]
Every TAR path from $S_{64}$ to an improvement consequently contains a set of
size at most $109$.  The compatible profile gives a strictly stronger barrier
than the individual profile on this instance.

\section{Certification and evidence levels}
\label{sec:certification}

Computer-assisted extremal results require two separate arguments: the CNF or
search instance must faithfully encode the mathematical problem, and the
claimed solver outcome must be checkable.  A second solver using the same
generator is useful testing but is not an independent certificate.  Certified
SAT workflows for combinatorial mathematics motivate our standard here
\cite{heule2016pythagorean,cruzfilipe2017rat,szeider2026lratcatcher}.

\subsection{Independently audited cell
\texorpdfstring{$u_{S_{100}}(6)=4$}{uS100(6)=4}}

The lower bound has an explicit witness
\begin{align*}
R={}&\{(26,59),(7,48),(73,59),(92,48),(96,2),(3,2)\},\\
\calU_{S_{100}}(R)={}&\{(4,1),(10,43),(89,43),(95,1)\}.
\end{align*}
An independent integer-geometry implementation verifies that these four and
only these four points are unlocked by $R$.

For the upper bound, a separate generator extracted the record from the public
notebook, checked its validity, rebuilt all conflicts by distance rings, and
constructed the decision instance ``six deletions unlock at least five outside
points.''  The grid has $9{,}836$ outsiders and $154{,}440$ link-graph edges in
total.  A simple certificate of seven pairwise disjoint conflict edges excludes
$9{,}099$ outsiders, because six deleted record vertices cannot cover such a
matching.  The independently checked matching file leaves $737$ conservative
candidates in the CNF.

The archived instance has $5{,}703$ variables and $24{,}775$ clauses.  Its
identity is fixed by the following artifact paths and hashes:
\begin{description}
\item[CNF]
\path{research/audit/r6_t5_matching.cnf}\\
SHA-256:
\path{f72229bacc773b19bf78209134cc666496854e60d1a515b4809ae2deeca23651}
\item[DRAT proof]
\path{research/audit/r6_t5_matching.drat}\\
SHA-256:
\path{149d3592fe5a2d6afccc05905b45ddc638947d537df6bc99865b9d190b20be9a}
\end{description}
The CNF is $381{,}497$ bytes.  CaDiCaL~1.9.5 at commit
\path{146207318796f094dcded87349a64f0c6927309e} produced a
$291{,}778{,}586$-byte binary DRAT proof.  \texttt{drat-trim} at commit
\path{2e3b2dc0ecf938addbd779d42877b6ed69d9a985} returned
\texttt{s VERIFIED} in $228.846$ seconds.  The proof manifest records all
commands, hashes, versions, and checker output.

This proof validates the archived CNF.  Confidence in the mathematical
translation comes from the independent generator, direct witness checks,
matching certificates, exhaustive small cardinality tests, and a written proof
of the encoding.  It is not a formal proof of the geometry-to-CNF generator.
A conversion of the full instance to LRAT followed by a CakeML-produced checker
has not been completed because of local storage constraints; only the checker
chain on a tiny smoke instance has passed.  The manuscript therefore does not
claim end-to-end formal verification.

\subsection{Machine-verified radius-seven exchange bound}

For $S_{100}$ at $r=7$, there are $1{,}324$ eligible outsiders.  Of the
$\binom{1324}{2}=875{,}826$ pairs, pair preprocessing rejects $863{,}884$ and
retains $11{,}942$.  An independently written pair-cache regenerator agrees
with those decisions and hashes.  The antichain DFS reports \textsc{unsat} for
six additions after $15{,}227$ nodes, $1{,}013{,}058$ family extensions, and
$475{,}262$ generated masks.  This is now corroborative rather than the sole
upper-bound evidence.

Any six-vertex clique lies in the $5$-core of the pair-compatibility graph.
The exact core has $392$ candidates and $11{,}054$ compatible edges.  A second
SAT encoding imposes the one-outsider cover clauses, every two-outsider forced
deletion mask, and all three-outsider incompatibilities.  Independent
standard-library programs regenerate the eligible candidates, pair graph,
$5$-core, and all $200{,}796$ ternary no-goods.  The resulting no-symmetry CNF
has $7{,}386$ variables and $308{,}590$ clauses:
\begin{description}
\item[Core CNF]
\path{research/audit/radius7_core_r7_k6_nosym.cnf}\\
SHA-256:
\path{2461ee72b797a636de8154422544c2c3673e43d9c83e069f41a1c434eef26bba}
\item[DRAT proof]
\path{research/audit/radius7_core_r7_k6_nosym.drat}\\
SHA-256:
\path{6a845c1af080a63661599f9633910b170293502c9319688cf55829ca1ba37c3b}
\end{description}
The CNF is $4{,}855{,}298$ bytes.  CaDiCaL returned \textsc{unsat} with exit
code $20$ and wrote a $25{,}282{,}853$-byte proof.  Independent
\texttt{drat-trim} verification returned \texttt{s VERIFIED} with exit code
$0$ in $26.669$ seconds; its core contained $252{,}110$ of the $308{,}590$
input clauses and used $22{,}191{,}309$ resolution steps, with zero RAT lemmas
in the core.  Together with the directly checked five-addition witness, this
establishes $e_{S_{100}}(7)=5$ at the machine-verified evidence level.

The DRAT checker proves that the archived CNF is unsatisfiable.  The independent
regenerators and the human proof of the reduction connect that CNF to the
geometric exchange problem, but this connection is not formally verified in a
proof assistant.  The full proof has not been converted to LRAT, and neither
the manuscript nor the artifacts have been peer reviewed.  ``Machine
verified'' is therefore the accurate description; ``formally verified
end-to-end'' is not.

\subsection{Independently audited radius-eight exhaustive search}

For $S_{100}$ at $r=8$, the search input has $2{,}036$ eligible outsiders and
$150{,}848$ inclusion-minimal one-outsider covers.  Its pair preprocessing
stores $331{,}300$ nonempty forced-deletion masks and retains $35{,}950$ of the
$2{,}071{,}630$ outsider pairs.  A standard-library audit, using an alternative
vertex-cover enumerator and direct squared-distance geometry, independently
regenerated the record, every candidate and cover, every forced mask, and
every pair decision.  All counts and decisions matched the binary input and
JSON cache exactly.

The positive side is the $8\to6$ witness displayed above, checked by a
separate exact-integer cubic enumeration of all triples in the resulting
$162$-point set.  For the upper bound, the C++ implementation of
Proposition~\ref{prop:antichain} exhaustively rejected a target of seven.  A
deliberately weaker replay removed every all-outsider-triple constraint and
disabled the coloring bound.  It still returned \textsc{unsat}, after
$78{,}031$ nodes, $10{,}009{,}190$ antichain extensions,
$5{,}388{,}204{,}289$ cover-union tests, and $3{,}217{,}876$ generated masks.
Every genuine seven-point exchange is feasible in this relaxation, so this
negative result yields $e_{S_{100}}(8)<7$ without relying on either the
all-outsider-triple code or color pruning.  Together with the witness, it gives
$e_{S_{100}}(8)=6$.

The source was independently recompiled and the decisive relaxed run replayed
with identical structural counters.  Regression tests reproduce the certified
radius-seven outcome; sanitizer and randomized differential tests exercise the
mask recurrence and coloring bound.  These checks materially narrow the trust
boundary, but they are not a portable proof trace.  No DRAT/LRAT certificate or
proof-assistant verification exists for the radius-eight upper bound.  We
therefore describe it as an exact exhaustive computation, independently
audited and cross-checked, rather than as formally or proof-trace verified.

The $S_{64}$ runs and the remaining table cells have executable witnesses and
current exact-solver logs, but do not yet have a uniform suite of independently
checked proof certificates comparable to the decisive $S_{100}$ radius-six and
radius-seven artifacts.  They are accordingly reported as reproducible
computations rather than proof-trace-certified results.

\section{Algorithmic interpretation}

The profiles have an immediate use in destroy-and-repair search.  After
deleting $r$ elements of an incumbent $S$, no feasible repair can insert more
than $e_S(r)$ outsiders, and the cheaper profile $u_S(r)$ is a safe upper
bound.  A radius satisfying $u_S(r)\le r$ cannot directly improve $S$.
Consequently, a solver can skip such radii when it seeks a one-shot
improvement, or deliberately treat them only as transit states in a
non-monotone search.

For $S_{100}$, the audited values exclude a direct improvement through radius
eight.  Any successful one-shot destroy-and-repair search must therefore
destroy at least nine incumbent points; the next record-improvement decision
is whether ten outsiders can be inserted after nine deletions.  In a global
model seeking at least $165$ points, let $x_s$ indicate whether the record
point $s$ is retained.  The same result gives the valid incumbent-overlap cut
\[
  \sum_{s\in S_{100}}x_s\le155.
\]
It excludes all $\sum_{r=0}^{8}\binom{164}{r}=11{,}494{,}182{,}836{,}236$
deletion masks in the first nine shells before any repair variables are
considered.  For
$S_{64}$, the difference between $u$ and $e$ at radius seven shows that ranking
a destroy set only by the number of individually unlocked points can be
misleading: the best unlocked collection need not contain a compatible repair
of the same size.

These observations are algorithmic consequences, not evidence that the
profiles predict heuristic performance in general.  Establishing such a link
would require multiple nonisomorphic constructions, repeated solver runs, and
an experimental protocol fixed before evaluating the outcomes.

\section{Reproducibility and provenance}
\label{sec:reproducibility}

The accompanying research directory is intended to be the reproducibility
package.  The key files are:
\begin{itemize}[leftmargin=*]
  \item \path{research/audit/AUDIT.md}, the independent radius-six audit;
  \item \path{research/audit/proof_manifest.json} and
  \path{research/audit/certificate_checksums.sha256}, recording the instance,
  proof, tools, and hashes;
  \item \path{research/audit/verify_certificate.sh}, the one-command replay of
  checksums, matching witnesses, and DRAT verification;
  \item \path{research/radius7_exchange_dfs.py} and
  \path{research/radius7_pair_compatibility.py}, the exchange solver and pair
  preprocessor;
  \item \path{research/radius7_dfs_n100_r7_k5.txt} and
  \path{research/radius7_dfs_n100_r7_k6.txt}, the radius-seven witness and
  current upper-bound log;
  \item \path{research/radius7_dfs_n64_r6_k5.txt},
  \path{research/radius7_dfs_n64_r7_k5.txt}, and
  \path{research/radius7_dfs_n64_r7_k6.txt}, the compatible $n=64$ logs;
  \item \path{research/audit/radius7_pair_cache_audit.json}, the independent
  pair-cache audit;
  \item \path{research/audit/radius7_core_audit.json}, the independent core and
  ternary-no-good regeneration;
  \item \path{research/audit/radius7_core_r7_k6_nosym.cnf} and
  \path{research/audit/radius7_core_r7_k6_nosym.drat}, the certified
  radius-seven upper-bound instance and proof;
  \item \path{research/radius8_cpp_antichain_dfs.cpp} and
  \path{research/radius8_cpp_REPORT.md}, the radius-eight solver and detailed
  computational report;
  \item \path{research/radius8_cpp_n100_r8_k6.json},
  \path{research/radius8_cpp_n100_r8_k6_verify.json}, and
  \path{research/radius8_cpp_n100_r8_k7_relaxed_nocolor.json}, the positive
  witness, direct check, and decisive relaxed upper-bound run;
  \item \path{research/audit/radius8_cpp_inputs_independent.py},
  \path{research/radius8_cpp_input_audit.json}, and
  \path{research/audit/RADIUS8_CPP_AUDIT.md}, the independent radius-eight
  input and implementation audits;
  \item \path{research/audit/verify_radius8_computation.sh}, the one-command
  checksum, witness, differential-test, rebuild, and decisive replay wrapper;
  \item \path{research/RESULTS_MANIFEST.json} and
  \path{research/audit/verify_frontier_bundle.sh}, the machine-readable
  claim/evidence snapshot and its fast cross-check wrapper;
  \item \path{research/RADIUS9_THEORY_NOTE.md}, the transition-barrier,
  incumbent-cut, and local TAR-component deductions;
  \item \path{research/audit/verify_tar_transition_theorem.py}, an exhaustive
  finite-model regression over all hereditary families on ground sets of size
  at most four (a check of edge cases, not a substitute for the proof);
  \item \path{research/radius9_FRONTIER.md},
  \path{research/radius9_cpp_n100_r9_k10.json}, and
  \path{research/radius9_alt_search_20260827.json}, the explicitly
  \textsc{unknown} radius-nine frontier and bounded-search diagnostics;
  \item \path{research/radius9_u7_witness.json},
  \path{research/radius9_u7_witness_verify.json}, and
  \path{research/audit/radius9_mask_direct_check.py}, the radius-nine
  individual-unlock witness and direct exact-integer checks.
\end{itemize}

From the repository root, the radius-six certificate and its matching
witnesses are replayed by
\begin{verbatim}
research/audit/verify_certificate.sh
\end{verbatim}
The exact radius-seven DFS can be rerun, subject to the documented Python
environment, by
\begin{verbatim}
python research/radius7_exchange_dfs.py \
  --n 100 --radius 7 --target 6 \
  --pair-cache research/radius7_paircompat_n100_r7.json \
  --seconds 300 \
  --output research/radius7_dfs_n100_r7_k6.replay.txt
\end{verbatim}
The independently checkable proof is replayed by
\begin{verbatim}
research/audit/tools/drat-trim/drat-trim \
  research/audit/radius7_core_r7_k6_nosym.cnf \
  research/audit/radius7_core_r7_k6_nosym.drat
\end{verbatim}
The required success line is \texttt{s VERIFIED}.  Replaying the DFS alone is
not a certificate; replaying the DRAT checker verifies the archived CNF, while
the reduction audit remains a separate part of the trust chain.

The radius-eight checks and decisive replay are run together by
\begin{verbatim}
research/audit/verify_radius8_computation.sh
\end{verbatim}
Passing \texttt{--full-input-audit} additionally regenerates every candidate,
cover, forced mask, and pair decision.  Equivalently, the decisive relaxation
alone can be rebuilt and replayed by
\begin{verbatim}
g++ -O3 -DNDEBUG -std=c++20 \
  -DR8CPP_RELAX_OUTSIDE_TRIPLES -DR8CPP_NO_COLOR_PRUNING \
  research/radius8_cpp_antichain_dfs.cpp -o /tmp/r8-relaxed
/tmp/r8-relaxed \
  research/radius7_pair_input_n100_r8.bin \
  research/radius7_paircompat_n100_r8.json \
  7 600 /tmp/r8-k7-relaxed.json
\end{verbatim}
The expected status is \textsc{unsat} with exit code $20$.  This is a
deterministic replay of the audited exhaustive program, not an independently
checkable UNSAT certificate.

The AlphaEvolve source was pinned at repository commit
\path{8f447457957deac61e28bf1676746f0753b3b2f8}.  The source notebook has
SHA-256
\path{653accd24369672f1cc757b203dae311f1489895a1cb8e6242d3308fe657891d},
and the canonical serialized $S_{100}$ coordinates have SHA-256
\path{ffe9583db47dbcc7e29e12a7d148fb1388827ac9917d1aefa7eb5fee89b4e7d6}.
The accompanying archive preserves the sources, environment specification,
witnesses, CNFs, proofs, checkers, manifests, and manuscript under immutable
hashes.

\section{Limitations and open problems}

The present study has several material limitations.
\begin{enumerate}[leftmargin=*]
  \item The exact values $u_{S_{100}}(7)$, $u_{S_{100}}(8)$, and both profiles
  at radius nine are open.  A timeout or a failed heuristic search is neither
  evidence of satisfiability nor a certified upper bound.
  \item The upper bound $e_{S_{100}}(7)<6$ has a verified DRAT proof, but the
  geometry-to-CNF reduction is audited software rather than an end-to-end
  theorem in a proof assistant.  No full LRAT/CakeML check has been performed.
  \item The upper bound $e_{S_{100}}(8)<7$ is an independently audited and
  replayed exhaustive computation, but it has no DRAT/LRAT proof trace.  It
  therefore retains the C++ implementation, compiler, and hardware as part of
  its trust base.
  \item Certification is uneven across the two tables.  The decisive
  $S_{100}$ radius-six and radius-seven bounds have verified DRAT chains; other
  upper-bound cells need similarly portable evidence.
  \item The formulas $D^u_S$ and $D^e_S$ lower-bound a true TAR barrier but do
  not generally determine it.  Computing exact max--min reconfiguration values
  is a separate problem.
  \item Two principal configurations are insufficient for a general empirical
  law.  Profiles for all SAT-optimal small grids and for multiple
  nonisomorphic large configurations would test whether rigidity correlates
  with size or search difficulty.
  \item No result here proves $C(64)=112$, $C(100)=164$, or any other global
  optimum.  The asymptotic bound in \cite{croot2026largesieve} does not close
  these finite cases.
  \item The novelty statement is based on a targeted search of primary sources.
  It should be checked against MathSciNet and zbMATH and discussed with domain
  experts before asserting priority.
\end{enumerate}

Promising next questions include producing a portable proof trace for the
radius-eight rejection, converting the existing proofs to LRAT, solving
$u_{S_{100}}(7)$ and $u_{S_{100}}(8)$, profiling small globally optimal
instances, comparing nonisomorphic sets of equal size, and deciding the next
record frontier $e_{S_{100}}(9)\ge10$.  The radius-nine counters indicate that
caching each complete deletion mask's unlocked-candidate bitset is a concrete
next optimization: it can replace repeated minimal-cover scans by bitset
intersections.  A more
theoretical direction is to characterize hypergraph classes for which the
profile-derived TAR bound is attained.

\section*{Use of computational tools}

Generative language-model tools, including OpenAI Codex, assisted computational
exploration, code development, literature search, and drafting.  The human
author is responsible for the manuscript, its claims, references, and
accompanying artifacts.

\section*{Acknowledgments}

The public configurations studied here originate in the AlphaEvolve repository
and associated paper \cite{georgiev2025exploration,alphaevolve59}.

\bibliographystyle{alpha}
\bibliography{references}

\end{document}
